\begin{enumerate}
\item El eDSL en TypeScript desarrollado en este trabajo reduce la dificultad percibida por usuarios nuevos en Nix: sus resultados son estadísticamente superiores en las pruebas cruzadas de Mann-Whitney-Wilcoxon para la experiencia de usuario general (Cuadro~\ref{cuadro-oux-nfs}) y para los aspectos generales y pragmáticos (Cuadro~\ref{cuadro-calidades}). Sin embargo, no mostró mejoras objetivas en el tiempo promedio de completitud (Nixlang = 34.4 minutos, Nixty = 41.7 minutos) ni en el porcentaje de acierto (Nixlang = 66.8~\%, Nixty = 72.3~\%).
\item Los estudiantes de Computación sin experiencia en Nix conciben un lenguaje de manejo de paquetes como una serie de instrucciones imperativas: todos lo describen con órdenes como instalar, actualizar o eliminar, y 9 de 10 representan los paquetes como objetos JSON, como es habitual en Python y JavaScript, los lenguajes que más dominaban. El lenguaje de Nix, Nixlang, en cambio, describe el resultado deseado mediante funciones, y es ahí, junto con los mensajes de error, donde más tropiezan: tras 1~h 30~min, solo llegaron a la tercera de siete actividades.
\item Es viable escribir configuraciones de Nix en TypeScript. Nixty, el eDSL desarrollado, funciona como una librería: el usuario describe su proyecto con objetos y funciones de TypeScript, y Nixty genera el archivo de configuración de Nix (\texttt{flake.nix}) equivalente. La mayoría de las utilidades de la biblioteca estándar de Nix se sustituyen por las del propio TypeScript. Su funcionamiento se verifica con 97 pruebas automatizadas.
\item En conjunto, Nixty reduce la barrera de entrada percibida, pero no la medida. Se observan tamaños de efecto grandes en la Experiencia de Usuario General ($r_g = 0.65$) y en nueve de las cualidades General, Pragmática y Hedónica (Cuadro~\ref{cuadro-calidades}). Dado que la barrera que reporta la comunidad puede ser en parte percibida~\cite{NixCommunitySurvey2024}, un eDSL en un lenguaje conocido constituye una vía plausible para facilitar la adopción de Nix.
\item La presente propuesta es todavía una implementación temprana. Queda pendiente evaluar su impacto con una muestra más amplia, someterla al criterio de la comunidad de Nix y estudiar su interacción con agentes de inteligencia artificial, en la línea de~\cite{Schwaighofer2026}. Además, Nixlang es solo una de las interfaces a través de las cuales el usuario interactúa con Nix. Según los deseos de la comunidad~\cite{2025CommunitySurvey}, también es necesario realizar estudios de usabilidad sobre su documentación y mensajes de error.
\end{enumerate}
